% diagrama_causal.tex — Esquema causal propio (Observación 5, revisión externa)
% No linealidad fisicoquímica del encalado -> limitaciones del PID ->
% mecanismos de mejora mediante AA/control avanzado.
% Se incluye vía \input{} dentro de un entorno figure en 03_introduction.tex.
% Layout por coordenadas absolutas (x,y) en cm, no posicionamiento relativo
% encadenado, para evitar colisiones entre columnas de distinta altura.
\begin{tikzpicture}[
  font=\sffamily\scriptsize,
  box/.style = {draw, rounded corners=2pt, align=center, text width=3.0cm,
                minimum height=1.6cm, inner sep=3pt, line width=0.6pt},
  causa/.style   = {box, fill=blue!6},
  efecto/.style  = {box, fill=red!6},
  mejora/.style  = {box, fill=green!6},
  flecha/.style  = {-{Stealth[length=2mm]}, line width=0.6pt},
  etiqueta/.style = {font=\sffamily\bfseries\scriptsize, text=black!65},
]

% ── Encabezados de columna ─────────────────────────────────────────────────────
\node[etiqueta, align=center, text width=2.6cm] at (1.6, 9.5)
  {No linealidad\\fisicoquímica\\del encalado};
\node[etiqueta, align=center, text width=3.0cm] at (6.7, 9.5)
  {Limitación del\\control PID};
\node[etiqueta, align=center, text width=2.8cm] at (11.8, 9.5)
  {Mecanismos de\\AA / control\\avanzado};

% ── Columna 1: causa raíz fisicoquímica (y = 7.5, 4.6, 1.5) ───────────────────
\node[causa] (sigmoidal) at (1.6, 7.5)   {Curva de titulación\\\textit{sigmoidal}:
  ganancia variable en\\varios órdenes de\\magnitud cerca de\\pH~$\approx$~7.0--7.2};
\node[causa] (variabilidad) at (1.6, 4.6) {Variabilidad del\\jugo crudo\\
  (variedad, madurez,\\campaña)};
\node[causa] (fouling) at (1.6, 1.5)   {\textit{Fouling} de\\electrodos y
  \\asimetría DCS~vs.\\laboratorio};

% ── Columna 2: nodo central PID (y=4.6) y sus dos consecuencias (y=1.1) ───────
\node[efecto] (pid)   at (6.7, 4.6) {PID sintonizado en\\un punto de
  operación\\pierde desempeño\\al alejarse de él};
\node[efecto, text width=3.2cm] (oscil) at (4.8, 1.0) {Oscilaciones\\sostenidas y
  \\sobredosificación\\de cal};
\node[efecto, text width=3.2cm] (costo) at (8.6, 1.0) {Costo económico\\doble:
  pH~$<$~6.8~$\to$\\inversión de sacarosa;\\pH~$>$~7.4~$\to$\\sacaratos de calcio};

% ── Columna 3: mecanismos de mejora (y = 7.5, 4.6, 1.5) ───────────────────────
\node[mejora] (fuzzy) at (11.8, 7.5)   {Control neuro-difuso\\orientado a la zona
  \\de máxima sensibilidad\\(punto de equivalencia)};
\node[mejora] (nmpc)  at (11.8, 4.6) {NMPC / caja gris con\\identificación de
  \\sistema explícita de\\la no linealidad};
\node[mejora] (soft)  at (11.8, 1.5) {Sensores virtuales /\\\textit{data-driven}
  \\robustos a \textit{fouling}\\y a baja frecuencia\\de muestreo};

% ── Flechas causales (columna 1 -> columna 2) ─────────────────────────────────
\draw[flecha] (sigmoidal) -- (pid);
\draw[flecha] (variabilidad) -- (pid);
\draw[flecha] (fouling) -- (pid);
\draw[flecha] (pid) -- (oscil);
\draw[flecha] (pid) -- (costo);

% ── Flechas causales (columna 1 -> columna 3) ─────────────────────────────────
\draw[flecha] (sigmoidal) -- (fuzzy);
\draw[flecha] (variabilidad) -- (nmpc);
\draw[flecha] (fouling) -- (soft);

% ── Flechas de mitigación (columna 3 -> columna 2, discontinuas) ──────────────
\draw[flecha, dashed, blue!60!black] (fuzzy) -- (pid.north east)
  node[midway, above, sloped, etiqueta, text=blue!60!black] {\scriptsize mitiga};
\draw[flecha, dashed, blue!60!black] (nmpc) -- (pid.east)
  node[midway, above, sloped, etiqueta, text=blue!60!black] {\scriptsize mitiga};
\draw[flecha, dashed, blue!60!black] (soft) -- (pid.south east)
  node[midway, below, sloped, etiqueta, text=blue!60!black] {\scriptsize mitiga};

\end{tikzpicture}
